% Cover letter for DSFE (AIMS Press) submission. Times 12pt, British English.
% Build:  pdflatex cover_letter
\documentclass[12pt,a4paper]{article}
\usepackage[T1]{fontenc}
\usepackage[utf8]{inputenc}
\usepackage{newtxtext,newtxmath}
\usepackage[british]{babel}
\usepackage[a4paper,margin=2.0cm]{geometry}
\usepackage{setspace}
\setstretch{1.0}
\usepackage{parskip}
\setlength{\parskip}{0.55em}
\usepackage[colorlinks=true,linkcolor=black,urlcolor=blue!60!black]{hyperref}
\pagestyle{empty}

\begin{document}

\noindent [Date]

\bigskip
\noindent To the Editors\\
\emph{Data Science in Finance and Economics} (DSFE), AIMS Press\\
Special Issue: Machine Learning in Economics and Finance

\bigskip
\noindent Dear Editors,

I am pleased to submit my manuscript, ``Diagnosing and mitigating the neutral-class
deficit in financial sentiment analysis: An explainability-guided approach on
multi-source investor text,'' for consideration in the special issue on Machine
Learning in Economics and Finance.

\textbf{Why this paper is a valuable addition.} Domain-adaptive models such as
FinBERT are the default tool for financial sentiment classification, and their
weakness on the neutral class is widely acknowledged, yet it is almost always
reported as an aggregate error rate and left unexplained. This manuscript makes that
weakness the object of study. Its contribution is methodological: it shows that the
neutral-class deficit is not diffuse noise but a single, lexically localised, and
repairable failure mode, and it demonstrates an explainability-guided loop that
closes the gap between diagnosing a model error and fixing it. Using a purpose-built
multi-source benchmark of retail, microblog and news text, and proper paired
significance testing (a McNemar test and stratified-bootstrap confidence intervals
that comparable applied studies frequently omit), the paper establishes a significant
advantage of FinBERT over a finance-lexicon baseline. It then uses LIME as a
diagnostic instrument to trace the model's low neutral recall to the systematic
misattribution of contextually neutral risk vocabulary, and shows that a small,
targeted fine-tuning intervention recovers neutral recall by 18.4 percentage points
without degrading performance on the polarised classes.

\textbf{Relation to previously published work.} The study builds on the financial
sentiment literature (Araci, 2019; Malo et al., 2014) and on recent surveys of
explainable AI in finance (Yeo et al., 2025), which identify the use of
interpretability as a feedback mechanism, rather than as a presentation or offline
audit layer, as under-explored. To my knowledge, this paper is a direct
demonstration of that loop on realistic multi-source financial text, and it
foregrounds statistical rigour (paired testing and bootstrap intervals) that is
often missing from applied sentiment comparisons. Limitations, including benchmark
scale and the single-seed nature of the fine-tuning result, are stated plainly in
the manuscript, and reusable code for a cross-validated extension and an external
benchmark comparison accompanies the submission.

\textbf{Reviewers.} I have not nominated specific reviewers and am happy to leave
referee selection to the editorial team. I have no opposed reviewers to declare.

\textbf{Declarations.} I confirm that this manuscript is original, has not been
published elsewhere, and is not under consideration for publication by another
journal. All authors (sole author) agree to publication in AIMS Press Open Access
format under the Creative Commons Attribution License. I understand the submission
may be screened by the CrossCheck originality-detection service. I confirm that
permission has been obtained for any reproduced material. Generative-AI use is
disclosed in the manuscript in the dedicated declaration section. The Article
Processing Charge is covered by Guangzhou University.

Thank you for your consideration.

\bigskip
\noindent Yours sincerely,

\medskip
\noindent Dominic Hubble\\
Department of Computer Science, Loughborough University\\
Email: dominichubble@gmail.com\\
Tel: +44 (0)1509 222222

\end{document}
